\chapter[线性代数基础：从 PCA 到自编码器]{线性代数基础：\\从 PCA 到自编码器}
\section*{概述}
\addcontentsline{toc}{section}{概述}

主成分分析（Principal Component Analysis, PCA）是本书一个自然且易于进入的起点，主要有两个原因：
\begin{itemize}
   \item \textbf{搭建通向概率模型与生成模型的桥梁。}
    作为一种广为人知的降维方法，PCA 为本书后续将要展开的若干思想提供了自然的过渡。PCA 的重构表述会自然引向自编码器，在自编码器中，压缩与重构直接从数据中学习。其概率扩展——概率主成分分析（Probabilistic PCA, PPCA）——又从生成建模的角度重新解释了这一熟悉的确定性方法。这样一条从代数投影，到重构，再到概率潜变量建模的路径，预示了现代框架（例如变分自编码器（VAE）和扩散模型）中所采用的技术；在这些更灵活、表达能力更强的框架中，相似的原理会再次出现。

    \item \textbf{重新回顾线性代数中的核心概念。}  
    PCA 能够自然地展示线性代数中的若干关键思想——投影、基变换、特征值与特征向量——这些概念支撑着本书后续许多结构。生成模型运行在高维空间中，数据通过矩阵进行表示与变换。协方差矩阵、方差分析、矩阵分解以及基变换等概念会反复出现，成为理解数据结构的重要工具。
\end{itemize}

基于这些原因，PCA 是本书非常理想的起点。

\section*{概念导图}
\addcontentsline{toc}{section}{概念导图}

本章将按照一条循序渐进的路径展开：从基本线性变换出发，进入 PCA，再由 PCA 过渡到自编码器。每一步的安排都旨在建立直觉、自然连接后续内容，并为后续章节中的概率生成模型做好准备。

\begin{itemize}

    \item[\(\rightarrow\)] \textbf{首先把矩阵理解为一种变换。}  
    我们不会把矩阵--向量乘法只看成机械的计算规则，而是把它理解为向量和坐标系如何被重新塑造。这个几何视角不仅为 PCA 提供必要直觉，也为后续生成模型中反复出现的“基于变换理解模型”的视角打下基础。

    \item[\(\rightarrow\)] \textbf{随后建立矩阵乘法的两个互补视角，并把它们与 PCA 联系起来。}  
    \emph{列组合}视角说明矩阵如何定义一组新的方向，并由这些方向组合得到输出；\emph{行--点积}视角则从投影与对齐的角度解释矩阵乘法。这两种理解共同预示了 PCA 的两个核心作用：改变基，以及通过投影测量变化程度。

    \item[\(\rightarrow\)] \textbf{接下来，把 PCA 表述为寻找最大方差方向的问题，并引入特征值和特征向量作为求解这一问题的关键工具。}  
    将数据投影到候选方向上并测量其离散程度之后，PCA 就转化为一个优化问题。由于协方差矩阵是对称矩阵，其特征向量构成一个标准正交基；在这一基下改写目标函数后，可以看到最大方差方向就是最大特征值对应的特征向量。

    \item[\(\rightarrow\)] \textbf{随后把这一论证扩展到完整的降维与重构。}  
    最大的若干特征值对应的特征向量定义了一组按重要性排序的正交主方向，每一个方向都尽可能捕获剩余方差。由此自然得到投影、重构以及解释方差比，从而同时从几何和实践层面理解 PCA 如何压缩数据。

    \item[\(\rightarrow\)] \textbf{然后进一步揭示 PCA 作为基变换的深层结构。}  
    对协方差矩阵进行对角化，可以理解为一种相似变换：它保持总方差不变，但把方差重新组织到与主方向对齐、彼此去相关的坐标轴上。这一视角也自然引出了更一般的矩阵分解，例如奇异值分解（SVD）。

    \item[\(\rightarrow\)] \textbf{最后，关注点从方差转向重构，从而建立 PCA 与自编码器之间的桥梁。}  
    PCA 可以被重新解释为一个最小化重构误差的编码器--解码器模型；随后可以看到，线性自编码器会恢复同一个最优子空间，而非线性自编码器则把这一思想从线性子空间扩展到弯曲流形。这条路径既揭示了确定性重构模型的能力，也揭示了其局限，并由此引出下一章中的概率潜变量模型。

\end{itemize}

下面的图示对本章路径进行了总结：从基本线性变换，经由 PCA 与重构，最终到达本书后续将使用的自编码器桥梁。

\begin{center}
\begin{tikzpicture}[
  >=Latex,
  every node/.style={font=\small},
  box/.style={draw,rounded corners=2pt,align=center,inner sep=5pt,
             text width=3.05cm,minimum height=1.15cm},
  flow/.style={->,thick}
]
\node[font=\normalsize\bfseries] at (6.0,2.75)
  {线性代数基础概览};

\node[box] (mat) at (1.7,0.95) {
  \textbf{把矩阵理解为变换}\\[3pt]
  {\scriptsize 基向量的像决定\\整个线性映射}
};
\node[box] (proj) at (6.0,0.95) {
  \textbf{内积与投影}\\[3pt]
  {\scriptsize 对齐程度 $\rightarrow$ 所选方向\\上的坐标}
};
\node[box] (eig) at (10.3,0.95) {
  \textbf{特征向量与特征值}\\[3pt]
  {\scriptsize 自然方向\\与缩放强度}
};
\draw[flow] (mat) -- (proj);
\draw[flow] (proj) -- (eig);

\node[box] (pca) at (10.3,-2.30) {
  \textbf{PCA}\\[3pt]
  {\scriptsize 最大方差\\主子空间}
};
\node[box] (basis) at (6.0,-2.30) {
  \textbf{基变换与 SVD}\\[3pt]
  {\scriptsize 对角化协方差\\并揭示结构}
};
\node[box] (ae) at (1.7,-2.30) {
  \textbf{重构桥梁}\\[3pt]
  {\scriptsize 自编码器\\流形\\概率}
};
\draw[flow] (eig) -- (pca);
\draw[flow] (pca) -- (basis);
\draw[flow] (basis) -- (ae);
\end{tikzpicture}
\end{center}


\medskip
\noindent
\section{作为线性变换的矩阵：基方向、投影与复合}

理解矩阵时，最好不要只把它看成一张数字表，而应把它看成作用于向量的变换的表示。矩阵定义了一个\emph{线性变换}：它以一种保持加法与数乘结构的方式移动向量。理解矩阵如何作用于向量，是一项基础能力；它不仅对 PCA 这类线性方法重要，对于由一系列线性与非线性变换构成的更复杂生成模型同样重要。

\subsection{矩阵对基方向的变换}

\paragraph{标准基与恒等变换。}

先从最简单的情形开始。在二维空间中，标准基（也称单位基）向量为
\[
\hat{e}_1 =
\begin{bmatrix}1 \\ 0\end{bmatrix},
\qquad
\hat{e}_2 =
\begin{bmatrix}0 \\ 1\end{bmatrix}.
\]
\(\mathbb{R}^2\) 中的任意向量，都能够唯一地写成这两个方向的线性组合。

将这两个基向量作为矩阵的两列，就得到单位矩阵
\[
I =
\begin{bmatrix}
1 & 0 \\
0 & 1
\end{bmatrix},
\]
它表示一个使所有向量保持不变的变换。从几何上看，恒等变换不会改变任何向量，因此方向、长度和夹角都保持不变。

\paragraph{示例：旋转 \(90^\circ\)。}

为了观察矩阵如何主动改变空间，考虑
\[
R =
\begin{bmatrix}
0 & -1 \\
1 & 0
\end{bmatrix}.
\]
该矩阵的两列告诉我们基向量如何被变换：
\[
\hat{e}_1 \mapsto
\begin{bmatrix}0 \\ 1\end{bmatrix},
\qquad
\hat{e}_2 \mapsto
\begin{bmatrix}-1 \\ 0\end{bmatrix}.
\]
每一个基方向都逆时针旋转了 \(90^\circ\)。由于所有向量都是由这些基向量构成的，因此同样的旋转会一致地作用于整个平面。

这一例子体现了线性代数中的一个核心原则：
\begin{quote}
\emph{一个线性变换完全由它把各个基向量映射到哪里所决定。}
\end{quote}

\paragraph{超越旋转：一般线性变换。}

矩阵并不一定表示旋转。它的各列不要求彼此正交，也不要求具有单位长度。

例如，考虑
\[
A =
\begin{bmatrix}
1 & 1 \\
0 & 1
\end{bmatrix}.
\]
这里，第一个基向量保持不变，而第二个基向量从
\(\begin{bmatrix}0 \\ 1\end{bmatrix}\) 被映射为 \(\begin{bmatrix}1 \\ 1\end{bmatrix}\)。这一变换对应一个剪切（shear）：一个坐标轴发生倾斜，而另一个保持不变。

尽管旋转与剪切在几何效果上不同，但它们的描述方式完全一致：指定各个基向量被映射到哪里。一旦这些映射确定，变换对任意其他向量的作用也就完全确定。

图~\ref{fig:ch1-matrix-transformations} 把这一基向量视角几何化：变换后的基向量决定了整个坐标网格如何被矩阵带到新的位置。

\begin{figure}[H]
\centering
\begin{minipage}[t]{0.4\textwidth}
\vspace{0pt}\centering
\includegraphics[width=\linewidth]{figures/ch01/panels/fig1_1a_rotation.pdf}
\end{minipage}\hfill
\begin{minipage}[t]{0.4\textwidth}
\vspace{0pt}\centering
\includegraphics[width=\linewidth]{figures/ch01/panels/fig1_1b_shear.pdf}
\end{minipage}
\caption{\textbf{基方向决定线性变换。}
浅色虚线表示原始坐标网格，实线表示变换后的网格。箭头 $A\hat e_1$ 与 $A\hat e_2$ 是变换后的基方向，也就是矩阵 $A$ 的两列。旋转保持方形网格，而剪切会使一族网格线发生倾斜；一旦变换后的基方向确定，其他任意向量的变换都可由线性组合得到。}
\label{fig:ch1-matrix-transformations}
\end{figure}

下面转向矩阵乘法，考察这些变换如何作用于任意向量。

\subsection{理解矩阵乘法的两个视角}

\paragraph{矩阵--向量乘法：列组合视角。}

设矩阵 \(A\) 的两列分别为 \(v_1\) 和 \(v_2\)。对于向量 \(x = \begin{bmatrix}a \\ b\end{bmatrix}\)，矩阵--向量乘法给出
\[
Ax = a\,v_1 + b\,v_2.
\]

这个表达式揭示了理解矩阵乘法的第一种基本方式。矩阵的各列定义了一组新的方向，而向量中的各个分量指定这些方向应如何组合。因此，矩阵乘向量可以理解为：用矩阵的列方向作加权组合来构造输出。

这一解释与向量本身的构造方式直接对应。任意向量都可以写成
\[
x = a\,\hat{e}_1 + b\,\hat{e}_2,
\]
即标准基向量的线性组合。矩阵--向量乘法所做的事情，只是把这些标准基方向替换成矩阵 \(A\) 的各列所编码的变换后坐标方向。

在一般情况下，同样有
\[
Ax = \sum_{j=1}^n x_j\,a_j,
\]
其中 \(a_j\) 表示 \(A\) 的第 \(j\) 列。

\paragraph{矩阵--向量乘法：逐行点积视角。}

理解矩阵--向量乘法还有第二种同样重要的方式。此时不再聚焦于矩阵的列，而是考察输出向量的各个分量。

在这一视角下，\(Ax\) 的第 \(i\) 个分量由 \(A\) 的第 \(i\) 行与输入向量做点积得到。这种“行乘列”的计算并不是任意规定的数值规则，而具有明确的几何含义。

\begin{tcolorbox}[
  colback=blue!6,
  colframe=blue!55!black,
  title={为什么点积能够编码投影},
  boxrule=0.8pt,
  arc=2mm,
  left=2mm,
  right=2mm,
  top=1mm,
  bottom=1mm
]

\textbf{余弦相似度。}
两个向量 \(v_1\) 和 \(v_2\) 的余弦相似度为
\[
\cos\theta = \frac{v_1^\top v_2}{\|v_1\|\,\|v_2\|},
\]
其中 \(\theta\) 是两者之间的夹角。

\textbf{标量投影。}
对上式变形，可以得到 \(v_1\) 在 \(v_2\) 方向上的标量投影：
\[
\text{proj}_{v_2}(v_1)
=
\frac{v_1^\top v_2}{\|v_2\|}.
\]
该量衡量 \(v_1\) 沿 \(v_2\) 方向所包含的分量大小。

\textbf{点积作为缩放后的投影。}
因此，点积可以写成
\[
v_1^\top v_2
=
\|v_2\| \times \text{（\(v_1\) 在 \(v_2\) 上的投影）}.
\]
如果 \(v_2\) 是单位向量，那么点积就等于标量投影。

\textbf{向量投影。}
将标量投影乘以 \(v_2\) 的单位方向，就得到相应的\emph{向量投影}：
\[
\mathrm{proj}_{v_2}(v_1)
=
\frac{v_1^\top v_2}{\|v_2\|^2}\, v_2.
\]
它给出了 \(v_1\) 沿 \(v_2\) 方向的分量的向量表示。

\end{tcolorbox}
上述讨论说明，逐行点积计算中出现的点积，在几何上可以理解为对投影或对齐程度的度量。从这一角度看，\(Ax\) 的每一个分量都记录了输入向量与 \(A\) 的某一个行方向之间的对齐程度。

\paragraph{推广到矩阵--矩阵乘法。}

同样的思想可以自然推广到矩阵--矩阵乘法。两个矩阵相乘，对应两个线性变换的复合。

乘积 \(AB\) 可以逐列理解为
\[
AB =
\big[\,A b_1 \;\; A b_2 \;\; \cdots \;\; A b_p\,\big],
\]
即 \(B\) 的每一列分别经过 \(A\) 的变换。等价地，当它作用于向量 \(v\) 时，乘积 \(AB\) 表示两个连续发生的变换：
\[
(AB)v = A(Bv).
\]
向量先经过 \(B\) 的变换，再由 \(A\) 继续变换。每一个阶段仍然都可以使用前面介绍的列组合视角和逐行点积视角来理解。

\medskip
\noindent
综合起来，这些视角构成了 PCA 的直接线性代数基础，并会在我们逐步建立更具表达能力的生成模型时反复出现：
\begin{itemize}
\item 矩阵的列描述一个变换如何重塑空间；
\item 列组合解释变换后的向量如何构造；
\item 逐行点积说明如何提取投影；
\item 矩阵--矩阵乘法则把多个变换串联起来。
\end{itemize}

\section{特征向量、特征值与矩阵的自然坐标系}
在深入 PCA 之前，我们再引入一个理解矩阵如何作用于空间的基本工具。上一节中，我们把矩阵解释为一种线性变换，它通过混合、拉伸与投影不同方向来重塑向量。为了更精确地描述这种作用，现在希望找到一些方向，使变换在这些方向上的行为尽可能简单。

特征向量与特征值恰好提供了这种语言。它们确定了矩阵的一些特殊方向：沿这些方向，变换不会把向量与其他方向混合，而只会把它映射为自身的一个标量倍数。

\paragraph{特征值与特征向量。}
给定方阵 \(A \in \mathbb{R}^{d \times d}\)，若存在非零向量 \(v \in \mathbb{R}^d\) 使得
\[
A v = \lambda v
\]
对某个标量 \(\lambda \in \mathbb{R}\) 成立，则称 \(v\) 为 \(A\) 的一个特征向量，\(\lambda\) 为对应的特征值。这个等式说明，当 \(A\) 作用于 \(v\) 时，输出仍落在 \(v\) 张成的同一条直线上；其长度被缩放 \(|\lambda|\) 倍，如果 \(\lambda<0\)，方向还会发生反转。

对于一般向量，乘以 \(A\) 可能造成旋转、拉伸、剪切或其他形式的方向改变。特征向量则标识了不会与其他方向混合的特殊方向：输出仍位于原来的直线上，仅相差一个标量因子 \(\lambda\)。因此，可以把特征向量理解为矩阵 \(A\) 所定义变换下的\emph{纯缩放方向}。

\paragraph{特征分解：在矩阵的自然坐标系中表示矩阵。}
如果矩阵 \(A \in \mathbb{R}^{d \times d}\) 有 \(d\) 个线性无关的特征向量，就可以把它们按列组成矩阵 \(V\)，并把对应的特征值放在对角矩阵 \(\Lambda\) 的对角线上。

具体地，若 \(v_1, v_2, \dots, v_d\) 是 \(A\) 的特征向量，对应特征值分别为 \(\lambda_1, \lambda_2, \dots, \lambda_d\)，则
\[
A [\, v_1 \; v_2 \; \cdots \; v_d \,]
=
[\, A v_1 \; A v_2 \; \cdots \; A v_d \,]
=
[\, \lambda_1 v_1 \; \lambda_2 v_2 \; \cdots \; \lambda_d v_d \,].
\]

用紧凑的矩阵形式可写为
\[
A V = V \Lambda,
\]
其中
\[
V = [\, v_1 \; v_2 \; \cdots \; v_d \,],
\qquad
\Lambda =
\begin{bmatrix}
\lambda_1 & 0 & \cdots & 0 \\
0 & \lambda_2 & \cdots & 0 \\
\vdots & \vdots & \ddots & \vdots \\
0 & 0 & \cdots & \lambda_d
\end{bmatrix}.
\]

对上式整理可得特征分解
\[
A = V \Lambda V^{-1}.
\]

\begin{tcolorbox}[
  colback=blue!6,
  colframe=blue!55!black,
  title={直觉：把特征分解理解为基变换},
  boxrule=0.8pt,
  arc=2mm,
  left=2mm,
  right=2mm,
  top=1mm,
  bottom=1mm
]

线性代数中的一个核心思想是：同一个向量可以在不同的坐标系中进行描述。特征分解通过引入一个\emph{新的基}把这一点显式化；在这个新基中，矩阵的作用会变得格外简单。

矩阵 \(V\) 的各列由 \(A\) 的特征向量组成。这些特征向量为整个空间定义了一组新的坐标轴，也就是特征向量基。在这一基下表示一个向量，就是用它沿这些优选方向的坐标来描述该向量。

\medskip
\textbf{改变坐标。}
给定在标准基下表示的向量 \(x\)，我们通过
\[
x_{\text{eig}} = V^{-1} x
\]
进行坐标变换。这个操作并不会改变向量本身，只会改变向量的\emph{表示方式}，把它转换为特征向量坐标。

\medskip
\textbf{特征基下的简单作用。}
一旦把向量表示到特征向量基中，\(A\) 的作用就变得非常简单：每一个坐标可以彼此独立地缩放，缩放因子就是对应的特征值，而这些缩放由对角矩阵 \(\Lambda\) 表示：
\[
x_{\text{eig}} \;\mapsto\; \Lambda x_{\text{eig}}.
\]

\medskip
\textbf{回到原始基。}
完成缩放后，再通过
\[
x' = V(\Lambda x_{\text{eig}})
\]
把结果转换回原始坐标系，从而在标准基下重构该向量。

\medskip
综合起来，分解 \(A = V \Lambda V^{-1}\) 可以读作
\[
x
\;\xrightarrow{\;V^{-1}\;}
\text{基变换}
\;\xrightarrow{\;\Lambda\;}
\text{各方向独立缩放}
\;\xrightarrow{\;V\;}
\text{回到原始基}.
\]

因此，特征分解说明：\emph{只要选择合适的坐标系}，一个矩阵的作用就可以化为简单的缩放。
\end{tcolorbox}

\paragraph{对称矩阵这一特殊情形。}
对于一般方阵，实数域上的特征分解可能不存在；也可能出现复特征值，或者没有足够多线性无关的特征向量。然而，当 \(A\) 是实对称矩阵（\(A=A^\top\)）时，情况会大大简化。

此时有若干重要性质成立：
\begin{itemize}
    \item \(A\) 的所有特征值都是实数。
    \item 存在一个\emph{正交}矩阵 \(U\)，其各列为 \(A\) 的特征向量，使得
    \[
    A = U \Lambda U^\top,
    \]
    其中 \(U\) 的各列均为单位长度且两两正交，并满足 \(U^\top U = I\)。
\end{itemize}

这种特殊结构意味着，对称矩阵总可以通过标准正交基变换进行对角化，因此坐标变换本身会保持长度和夹角不变。正如后面将看到的，这一性质对 PCA 至关重要，因为经验协方差矩阵始终是对称矩阵。

至此，我们已经具备了表述并求解 PCA 问题所需的工具。

\section{通过方差最大化理解主成分分析}

主成分分析（PCA）\cite{pearson1901,hotelling1933} 通过引入一个适应数据结构的新坐标系，为数据降维提供了一套系统方法。其核心思想，是识别一组新的坐标轴，使数据沿这些轴的变化程度最大。将坐标轴与数据最为分散的方向对齐后，PCA 就揭示了数据在哪些方向上包含最多信息，并使我们能够在尽量少损失结构的前提下获得低维表示。

PCA 的直觉可以概括为两个核心原则：
\begin{enumerate}
    \item \textbf{把方差视为信息量。}  
    数据在某个方向上的方差越大，数据点通常越容易被区分，因此该方向携带的结构信息也越丰富。相反，方差很小的方向对区分数据的贡献较少。

    \item \textbf{顺序进行正交最大化。}  
    PCA 首先找出能够捕获最大可能方差的单一方向。随后继续选择新的方向，在保持与已选方向正交的约束下，尽可能捕获剩余方差。
\end{enumerate}

\subsection{PCA 的数学表述}

\paragraph{数据中心化与投影方差的定义。}
为了便于表述 PCA，首先把数据中心化，使整个数据云以原点为中心。给定数据点 \(\{x^{(i)}\}_{i=1}^n \subset \mathbb{R}^d\)，假设
\[
\frac{1}{n}\sum_{i=1}^n x^{(i)} = 0.
\]
中心化可以确保我们测量的是数据相对于其中心的变化程度，而不会受到任意平移偏置的影响。

前面在介绍投影和点积时已经讨论过：沿某个方向的变化程度，可以通过考察数据在该方向上的投影强度来理解。设 \(v \in \mathbb{R}^d\) 是一个单位向量，表示我们希望把数据投影到的候选方向，因此 \(\|v\|=1\)。对于数据点 \(x^{(i)}\)，其在 \(v\) 上的标量投影为
\[
v^\top x^{(i)},
\]
因为 \(v\) 的长度为 1。相应的向量投影为
\[
(v^\top x^{(i)})\,v.
\]
这再次说明，当 \(v\) 是单位向量时，点积恰好给出了沿该方向的标量投影。

由于数据已经中心化，投影标量值的均值同样为 0：
\[
\frac{1}{n}\sum_{i=1}^n v^\top x^{(i)}
=
v^\top \left(\frac{1}{n}\sum_{i=1}^n x^{(i)}\right)
=
0.
\]
为了量化这些投影值的离散程度，我们计算它们相对于该均值的方差。因此，数据沿方向 \(v\) 投影后的方差为
\[
\mathrm{Var}(v^\top X)
=
\frac{1}{n}\sum_{i=1}^n \bigl(v^\top x^{(i)}\bigr)^2.
\]
从几何上看，投影越大、分布越分散的方向，对应的方差越高。

\paragraph{协方差矩阵形式。}
投影方差可以用矩阵形式紧凑地写为
\[
\mathrm{Var}(v^\top X)
=
\frac{1}{n}\sum_{i=1}^n v^\top x^{(i)} (x^{(i)})^\top v
=
v^\top \left( \frac{1}{n}\sum_{i=1}^n x^{(i)} (x^{(i)})^\top \right) v.
\]
括号内的矩阵就是经验协方差矩阵
\[
\Sigma = \frac{1}{n} \sum_{i=1}^n x^{(i)} (x^{(i)})^\top.
\]
因此，数据沿方向 \(v\) 投影后的方差可以写成
\[
\mathrm{Var}(v^\top X) = v^\top \Sigma\, v.
\]

于是，PCA 问题就是在单位长度约束下寻找使上述量最大的方向：
\[
\max_{v \in \mathbb{R}^d} \; v^\top \Sigma v
\quad \text{约束为} \quad v^\top v = 1.
\]

\subsection{求解 PCA 优化问题}

协方差矩阵 \(\Sigma\) 按定义就是对称矩阵。如上一节所述，对称矩阵具有正交特征分解：
\[
\Sigma = U \Lambda U^\top,
\qquad
\Lambda = \mathrm{diag}(\lambda_1,\ldots,\lambda_d),
\]
其中 \(U\) 的各列是标准正交特征向量，特征值按
\[
\lambda_1 \ge \lambda_2 \ge \cdots \ge \lambda_d \ge 0
\]
排序。

为了求解优化问题，我们把目标函数表示到这一特征向量坐标系中。

\paragraph{步骤 1：变换到特征向量坐标。}
按照前面介绍的基变换解释，把方向向量 \(v\) 在特征向量基下表示为
\[
\tilde{v} = U^\top v.
\]
由于 \(U\) 为正交矩阵，\(U^\top = U^{-1}\)，因此该操作只是把 \(v\) 的坐标从标准基转换到特征向量基。

正交变换还会保持向量长度：
\[
\tilde{v}^\top \tilde{v}
=
(U^\top v)^\top (U^\top v)
=
v^\top U U^\top v
=
v^\top v.
\]
因此，单位长度约束 \(v^\top v = 1\) 会转化为
\[
\tilde{v}^\top \tilde{v} = 1.
\]

\paragraph{步骤 2：把目标函数对角化。}
将特征分解代入目标函数，有
\[
v^\top \Sigma v
=
v^\top U \Lambda U^\top v
=
(U^\top v)^\top \Lambda (U^\top v)
=
\tilde{v}^\top \Lambda \tilde{v}.
\]
由于 \(\Lambda\) 为对角矩阵，上式化简为
\[
\tilde{v}^\top \Lambda \tilde{v}
=
\sum_{i=1}^d \lambda_i \tilde{v}_i^2.
\]
这明确表明：沿方向 \(v\) 捕获到的方差，是各个特征值的加权和，而权重由 \(v\) 与各特征向量的对齐程度决定。

\paragraph{步骤 3：给出方差的上界。}
由于 \(\tilde{v}_i^2 \ge 0\)，且 \(\sum_{i=1}^d \tilde{v}_i^2 = 1\)，目标函数实际上是各特征值的加权平均：
\[
\sum_{i=1}^d \lambda_i \tilde{v}_i^2
\le
\lambda_1 \sum_{i=1}^d \tilde{v}_i^2
=
\lambda_1.
\]
任何较小特征值的组合都不可能超过最大的特征值。

\paragraph{步骤 4：得到最优方向。}
只有当全部权重都集中到最大特征值上时，才会达到该上界：
\[
\tilde{v} = [1,\,0,\,\ldots,\,0]^\top.
\]
将其变换回原始坐标，得到
\[
v = U \tilde{v} = u_1,
\]
即最大特征值 \(\lambda_1\) 对应的特征向量。\footnote{译注：这里给出的是一个最优解。若最大特征值为单重，则正、负两个单位特征向量都达到上界；若最大特征值为重根，则其特征子空间中的任意单位向量都达到上界。原文的坐标选择不表示最优解唯一。}

因此，使投影方差最大的方向，恰好就是协方差矩阵的第一特征向量，而最大方差等于 \(\lambda_1\)。

\paragraph{后续主成分。}
因此，第一个主成分就是 \(u_1\)。其余主成分通过重复相同的优化过程获得，同时要求新方向与此前已经选择的方向正交。第 \(k\) 个主成分就是第 \(k\) 大特征值对应的特征向量 \(u_k\)，并满足
\[
u_k^\top u_j = 0 \qquad \text{对所有 } j < k \text{ 成立}.
\]

从几何上看，每一个新的主方向都会在保持与此前方向正交的同时，尽可能捕获尚未解释的变化。由于这些方向是协方差矩阵的正交特征向量，由它们得到的主成分坐标彼此不相关。协方差矩阵的特征分解自然地给出了这样一组有序、标准正交的方向。

\subsection{用于降维的 PCA}

一旦确定主成分，PCA 就提供了一套具体的降维过程：先把数据投影到一个低维子空间中；如果需要，还可以再把它重构回原始空间。

设
\[
U_k = [\,u_1\;u_2\;\cdots\;u_k\,] \in \mathbb{R}^{d \times k}
\]
表示由前 \(k\) 个主成分组成的矩阵，也就是协方差矩阵 \(\Sigma\) 中对应最大 \(k\) 个特征值 \(\lambda_1 \ge \lambda_2 \ge \cdots \ge \lambda_k\) 的特征向量。PCA 通过只使用这 \(k\) 个捕获最大方差的方向来表示数据点 \(x \in \mathbb{R}^d\)，从而降低其维度。

\paragraph{投影（编码）。}
\(x\) 的低维表示通过将其投影到主成分子空间得到：
\[
z = U_k^\top x \in \mathbb{R}^k.
\]
向量 \(z\) 包含 \(x\) 沿各主方向的坐标，构成数据的压缩表示。

\paragraph{重构（解码）。}
只使用这 \(k\) 个成分对 \(x\) 进行近似，可通过把低维表示映射回原始空间实现：
\[
\hat{x} = U_k z = U_k U_k^\top x.
\]
这一操作取得 \(x\) 在主成分空间中的坐标，再利用 \(U_k\) 中的主方向进行组合，从而在原始特征空间中重构出一个向量。由于只保留前 \(k\) 个主成分，因此重构过程中不可避免地会损失部分信息。

值得注意的是，当 \(k=d\) 时，有 \(U_k \in \mathbb{R}^{d\times d}\)，并且
\[
U_k U_k^\top = I,
\]
因此重构是精确的。而当 \(k<d\) 时，\(U_k U_k^\top\) 不再是单位矩阵，\(\hat{x}\) 只能是原始数据点的近似。

\paragraph{解释方差比。}
在实践中，需要量化降维后保留了原始总变化量中的多少。最常用的指标就是\emph{解释方差比}（explained variance ratio）。

原始数据集中的总方差为
\[
\sum_{i=1}^d \lambda_i,
\]
而前 \(k\) 个主成分捕获的方差为
\[
\sum_{i=1}^k \lambda_i.
\]
因此，\(k\) 维表示所保留的方差比例为
\[
\text{解释方差比}
=
\frac{\sum_{i=1}^k \lambda_i}{\sum_{i=1}^d \lambda_i},
\]
它为选择主成分数量提供了一个有理论依据、由数据驱动的标准，在实际应用中被广泛采用。

下面进一步仔细考察 PCA 的结构。这将说明为什么解释方差比能够以一种有原则的方式量化降维所保留的信息量，并揭示下一节将讨论的 PCA 更深层本质。

\section{把 PCA 理解为基变换：相似变换与 SVD}
现在从更深的结构视角重新审视 PCA。上一节中，PCA 是作为方差最大化问题的解得到的。本节将说明，同一过程也可以理解为一种\emph{基变换}：它将协方差矩阵对角化，并把方差重新组织到一组正交方向上。这一视角能够解释为什么总方差保持不变、为什么特征值能够自然量化各主方向上捕获的方差，以及为什么 PCA 会产生彼此不相关的坐标。

这一视角还会自然引出奇异值分解（Singular Value Decomposition, SVD）。SVD 将同样的几何思想推广到协方差矩阵之外，也推广到非方阵和非对称矩阵。虽然基本形式的 PCA 并不要求掌握 SVD，但它提供了更一般的线性代数视角，并将在本书后续内容中持续发挥作用。

为了建立这种解释，先介绍两个关键概念：矩阵的迹以及相似变换。

\subsection{迹、相似变换与方差保持}

\paragraph{协方差矩阵的迹。}
对于任意方阵 \(A\in\mathbb{R}^{d\times d}\)，其迹定义为对角元素之和：
\[
\mathrm{Tr}(A)=\sum_{i=1}^d a_{ii}.
\]
当 \(A\) 是协方差矩阵 \(\Sigma\) 时，每个对角元素 \(\sigma_{ii}\) 都是第 \(i\) 个坐标的方差。因此
\[
\mathrm{Tr}(\Sigma)
=
\sum_{i=1}^d \Sigma_{ii}
=
\sum_{i=1}^d \mathrm{Var}(X_i),
\]
这恰好就是数据集沿所有坐标方向的总方差。

这一点对 PCA 很重要。基变换虽然可能改变方差在各个坐标上的分配方式，但总方差不会改变。\footnote{译注：这里关于总方差的结论应限于 PCA 所用的标准正交基变换。一般坐标变换 \(y=A^{-1}x\) 给出 \(\operatorname{Cov}(y)=A^{-1}\Sigma A^{-\top}\)，并非一般的相似变换，总方差通常会改变。对线性算子而言，相似变换保持迹的代数结论仍成立。}矩阵的迹为表达这一总方差提供了紧凑方式。

\paragraph{相似变换。}
前面已经看到，若一个新基由可逆矩阵 \(A\) 定义，则把向量表示到该基中对应坐标变换
\[
v_A = A^{-1}v.
\]
同样的思想也适用于线性变换本身。当线性变换 \(T\) 在这个新基中进行表示时，其矩阵表示变为
\[
T_A = A^{-1}TA,
\]
这称为\emph{相似变换}（similarity transformation）。它表示的仍然是\emph{同一个}线性变换，只是改用由基 \(A\) 决定的另一套坐标系统进行描述，如下推导所示。

\begin{tcolorbox}[
  colback=blue!6,
  colframe=blue!55!black,
  title={为什么 $T_A = A^{-1} T A$},
  boxrule=0.8pt,
  arc=2mm,
  left=2mm,
  right=2mm,
  top=1mm,
  bottom=1mm
]
若 \(v_A=A^{-1}v\) 表示向量 \(v\) 在由 \(A\) 定义的新基中的坐标，那么线性变换 \(T\) 在这个坐标系中也必须具有对应表示 \(T_A\)。两个矩阵描述的是\emph{同一个}几何变换，因此，在新坐标中应用 \(T_A\)，必须与先在原坐标中应用 \(T\)、再改变基得到相同结果：
\[
v'_A = T_A v_A,
\qquad
v'_A = A^{-1}(Tv).
\]
由于这两条路径必须在 \(A\) 坐标中产生同一个向量，代入 \(v_A=A^{-1}v\) 得
\[
T_A A^{-1}v = A^{-1}Tv
\quad\text{对所有 }v\text{ 成立}.
\]
两边左乘 \(A\)，有
\[
A T_A A^{-1}=T,
\qquad\Longrightarrow\qquad
T_A=A^{-1}TA.
\]
\end{tcolorbox}

\paragraph{基变换下迹的不变性。}

迹的一个关键代数性质是循环恒等式
\[
\mathrm{Tr}(BC)=\mathrm{Tr}(CB),
\]
只要矩阵乘积 \(BC\) 与 \(CB\) 均有定义即可。

结合这一性质与相似变换 \(T_A=A^{-1}TA\)，可得
\[
\mathrm{Tr}(T_A)
=
\mathrm{Tr}(A^{-1}TA)
=
\mathrm{Tr}(T).
\]

这说明迹在相似变换下保持不变。因此，虽然基变换可能重新分配方差在各个坐标上的表达方式，但不会改变总方差。

\subsection{把 PCA 理解为相似变换}

现在从基变换视角重新审视 PCA。对于单个数据点 \(x\in\mathbb{R}^d\)，低维表示通过投影到主成分方向获得，即 \(z=U_k^\top x\)。在整个数据集层面，如果把中心化后的数据点按行收集到矩阵 \(X\) 中，那么 PCA 对应线性变换 \(Z=XU\)，其中 \(U\) 是正交矩阵，其各列由协方差矩阵 \(\Sigma\) 的特征向量组成。

\medskip
\noindent
下面考察这种基变换如何改变数据的\emph{协方差结构}。

\paragraph{变换后数据的协方差。}
变换后数据的协方差矩阵为
\[
\tilde{\Sigma} = \mathrm{Cov}(Z)=\mathrm{Cov}(XU).
\]

一般地，随机向量 \(Y\) 的协方差可以写为
\[
\mathrm{Cov}(Y)
=\mathbb{E}\big[(Y-\mathbb{E}[Y])^\top(Y-\mathbb{E}[Y])\big].
\]
在当前情形下，数据已经中心化，因此 \(\mathbb{E}[X]=0\)。由于 \(U\) 是固定的线性变换，有
\[
\mathbb{E}[Z]=\mathbb{E}[XU]=\mathbb{E}[X]U=0.
\]
于是变换后数据的协方差化简为
\[
\tilde{\Sigma}
=\mathbb{E}\big[(Z-\mathbb{E}[Z])^\top(Z-\mathbb{E}[Z])\big]
=\mathbb{E}[Z^\top Z].
\]

代入 \(Z=XU\) 得
\[
\tilde{\Sigma}
=\mathbb{E}[(XU)^\top(XU)]
=U^\top\mathbb{E}[X^\top X]U
=U^\top\Sigma U.
\]

\medskip
\noindent
\textbf{PCA 之前与之后。}
现在可以明确写出这一核心结构结果：
\[
\mathrm{Cov}(X)=\Sigma,
\qquad
\mathrm{Cov}(Z)=U^\top\Sigma U.
\]
也就是说，PCA 按照
\[
\Sigma\;\longmapsto\;U^\top\Sigma U
\]
来变换协方差矩阵，这正是原协方差矩阵的一次\emph{相似变换}。

\paragraph{相似变换的结果。}
由于相似变换保持迹不变，有
\[
\mathrm{Tr}(\Sigma)
=
\mathrm{Tr}(U^\top\Sigma U)
=
\mathrm{Tr}(\Lambda),
\]
其中 \(\Sigma=U\Lambda U^\top\) 是协方差矩阵的特征分解。

这正体现了 PCA 的几何意义。PCA 之前，协方差矩阵 \(\Sigma\) 可能包含非对角项，表示不同坐标之间存在相关性。变换到特征向量基之后，协方差矩阵变成对角形式：
\[
U^\top\Sigma U
=
\Lambda
=
\mathrm{diag}(\lambda_1,\ldots,\lambda_d).
\]
因此，变换后的坐标彼此不相关，而每个主方向上的方差会以对应特征值的形式直接出现在对角线上。

PCA 并不会创造或消除方差。它只是把同样的总方差重新表达在一套与数据自然方向对齐的坐标系中。由迹表示的总方差保持不变，但其分布从相关坐标重新组织到了彼此正交的主方向上。

这也解释了为什么解释方差比具有如此自然的形式。保留前 \(k\) 个主成分，就保留方差 \(\sum_{i=1}^k\lambda_i\)，而总方差为 \(\sum_{i=1}^d\lambda_i\)。因此
\(\frac{\sum_{i=1}^k\lambda_i}{\sum_{i=1}^d\lambda_i}\)
度量的正是舍弃其余
方向之后所保留的总方差比例。

图~\ref{fig:ch1-pca-change-of-basis} 直接展示了这种坐标变换：同一数据云在主成分坐标中变得与坐标轴对齐，同时总方差保持不变。

\begin{figure}[H]
\centering
\begin{minipage}[t]{0.4\textwidth}
\vspace{0pt}\centering
\includegraphics[width=\linewidth]{figures/ch01/panels/fig1_3a_original_coordinates.pdf}
\end{minipage}\hfill
\begin{minipage}[t]{0.42\textwidth}
\vspace{0pt}\centering
\includegraphics[width=\linewidth]{figures/ch01/panels/fig1_3b_pca_coordinates.pdf}
\end{minipage}
\caption{\textbf{PCA 通过改变基来对角化协方差。}
图中分别在原始坐标和主成分坐标下展示同一组数据。在新基中，主方向成为坐标轴，各坐标之间的协方差消失，而总方差保持不变。}
\label{fig:ch1-pca-change-of-basis}
\end{figure}

\subsection{从特征分解到奇异值分解}

前面的讨论表明，PCA 可以理解为一种基变换：它把协方差矩阵对角化，并将方差重新组织到彼此正交的方向上。由此自然引出一个更一般的问题：
\begin{quote}
如果我们想理解的矩阵不是协方差矩阵，不是对称矩阵，甚至不是方阵，该怎么办？
\end{quote}
奇异值分解（Singular Value Decomposition, SVD）提供了相应的一般框架~\cite{strang2016introduction}。

特征分解建立在关系 \(Av=\lambda v\) 上，它寻找那些在矩阵作用后仍保持原方向、仅发生缩放的方向。沿这些方向，变换的行为尤其简单：输出仍处于同一条直线上，只被缩放 \(\lambda\) 倍。

这一视角很强大，但也存在限制。它只适用于方阵，并且通常还需要额外条件才能保证存在完整且性质良好的实数分解。

SVD 对这一思想进行了推广。它不再要求寻找“被矩阵作用后仍保持自身方向”的向量，而提出一个更灵活的问题：\emph{哪些相互正交的输入方向会被映射到哪些相互正交的输出方向？这些方向分别被拉伸多少？}

\paragraph{SVD 视角。}
对于实矩阵 \(A\in\mathbb{R}^{m\times n}\)，奇异值分解写为
\[
A = U S V^\top,
\]
其中，\(V\in\mathbb{R}^{n\times n}\) 包含一组标准正交的输入方向；\(U\in\mathbb{R}^{m\times m}\) 包含一组标准正交的输出方向；而 \(S\in\mathbb{R}^{m\times n}\) 是广义意义上的矩形对角矩阵，其主对角线上为非负数
\(\sigma_1\ge\sigma_2\ge\cdots\ge0\)。

这些量满足
\[
A v_i = \sigma_i u_i,
\qquad
A^\top u_i = \sigma_i v_i,
\]
其中 \(v_i\) 是右奇异向量，\(u_i\) 是对应的左奇异向量，\(\sigma_i\) 是对应的奇异值。

从几何上看，SVD 表明每个正交输入方向 \(v_i\) 会被映射为某个正交输出方向 \(u_i\)，并由奇异值 \(\sigma_i\) 决定缩放倍数。与特征分解不同，输入方向和输出方向不需要重合；这也正是为什么 SVD 对\emph{任意}实矩阵都存在，包括矩形矩阵和非对称矩阵。于是，特征分解可以看作一种特殊情形：当 \(A\) 为对称矩阵时，输入和输出方向重合（\(U=V\)），奇异值等于特征值的绝对值。\footnote{译注：原文此处的 \(U=V\) 还要求 \(A\) 为半正定矩阵。对于含负特征值的实对称矩阵，奇异值仍等于特征值的绝对值，但负特征值对应的左右奇异向量方向相反。例如 \(A=[-1]\) 的奇异值为 1，左右奇异向量不能同时取同一个单位向量。}

SVD 再次强化了本书中一个会反复出现的核心主题：一旦选择了合适的坐标系统或分解方式，复杂的线性变换往往会变得容易理解得多。PCA 对协方差矩阵所做的事情，就是旋转到特征向量基中；SVD 则把同样的思想推广到任意矩阵。

即使后续模型超越了线性变换，这种结构化的思考习惯仍然很重要。我们会不断寻找新的表示，使复杂运算能够分解成更简单、更容易解释的组成部分。从这一意义上说，SVD 不只是一个数值工具，更是 PCA 背后同一套几何思想的概念推广。

\begin{tcolorbox}[
  colback=blue!6,
  colframe=blue!55!black,
  title={秩一结构与矩阵的几何含义},
  boxrule=0.8pt,
  arc=2mm,
  left=2mm,
  right=2mm,
  top=1mm,
  bottom=1mm
]

SVD 最具启发性的解释之一，是它的秩一展开：
\[
A = \sum_{i=1}^r \sigma_i\,u_i v_i^\top,
\]
其中矩阵的秩 \(r\) 表示构造该矩阵所需的独立方向数量。

每一项 \(\sigma_i u_i v_i^\top\) 都是一个\emph{秩一变换}，由两个向量的外积构成。从几何上看，它执行一个简单的三步过程：
\[
\text{投影到 }v_i
\;\rightarrow\;
\text{按 }\sigma_i\text{ 缩放}
\;\rightarrow\;
\text{沿 }u_i\text{ 重新定向}.
\]

由于 \(v_i^\top x\) 是一个标量，输出 \(u_i v_i^\top x\) 总是单一向量 \(u_i\) 的某个标量倍数。因此所有可能输出都位于同一个固定方向上，该变换最多只能张成一维输出空间，所以秩为 1。这说明任意矩阵都可以理解为这些基本外积变换的叠加，其中每一项都捕获一种独立的作用模式。

\end{tcolorbox}

\section{从 PCA 到自编码器：重构、子空间与流形}

前面几节主要从\emph{方差最大化}的角度建立 PCA：我们寻找数据最为分散的方向，并把这些方向作为新的坐标系。这个视角揭示了 PCA 如何重新组织数据，使最有信息量的变化显式呈现出来。

不过，还有另一种同样自然且非常直观的方式来理解同一个方法。我们可以不再问“哪些方向捕获了最大的方差”，而改问一个重构问题：
\begin{quote}
\emph{如果先把数据压缩为一个 \(k\) 维表示，再从中进行重构，那么哪一种 \(k\) 维表示损失的信息最少？}
\end{quote}
这种重新表述把 PCA 从一个特征值问题转变为一个近似问题。几何结构并没有改变，但关注重点由\emph{捕获方差}转向了\emph{最小化重构误差}。这个转变很重要，因为在自编码器以及许多后续表示学习模型中，真正起组织作用的原则将不再是方差最大化，而是重构。

\subsection{从方差最大化到重构}

设 \(U_k=[\,u_1,\dots,u_k\,]\in\mathbb{R}^{d\times k}\) 表示由协方差矩阵前 \(k\) 个特征向量组成的矩阵。对于中心化数据点 \(x_i\in\mathbb{R}^d\)，PCA 给出低维表示
\[
z_i = U_k^\top x_i,
\]
它记录了 \(x_i\) 沿各主方向的坐标。从该表示进行重构得到
\[
\hat{x}_i = U_k z_i = U_k U_k^\top x_i,
\]
这恰好是 \(x_i\) 到由 \(U_k\) 各列张成的 \(k\) 维子空间上的正交投影。

有了原始点 \(x_i\) 和其重构 \(\hat{x}_i\) 后，信息损失最自然的度量就是两者之差。由此得到平方重构误差
\[
\|x_i-\hat{x}_i\|^2
=
\|x_i-U_kU_k^\top x_i\|^2.
\]
对所有数据点求和可得总重构损失
\[
L(U_k)
=
\sum_{i=1}^n \|x_i-U_kU_k^\top x_i\|^2.
\]

从这一视角看，PCA 可以写成
\[
\min_{U_k^\top U_k=I_k}
\;
\sum_{i=1}^n \|x_i-U_kU_k^\top x_i\|^2,
\]
它所回答的问题是：
\begin{quote}
\emph{哪一个具有标准正交基的 \(k\) 维子空间，能够使数据压缩后再重构所造成的失真最小？}
\end{quote}

这一重构形式与前面的 PCA 方差最大化形式完全等价，并与经典的最佳低秩近似视角密切相关~\cite{jolliffe2016pca}。虽然目标函数写法不同，但两者都会得到同一个最优 \(k\) 维子空间，即由协方差矩阵最大的若干特征值对应的特征向量所张成的子空间。

图~\ref{fig:ch1-pca-variance-geometry} 在二维情形下展示了二者的等价性：主方向既使投影后的坐标最为分散，也使正交重构残差最小。

\begin{figure}[H]
\centering
\begin{minipage}[t]{0.42\textwidth}
\vspace{0pt}\centering
\includegraphics[width=\linewidth]{figures/ch01/panels/fig1_2a_bad_direction.pdf}
\end{minipage}\hfill
\begin{minipage}[t]{0.42\textwidth}
\vspace{0pt}\centering
\includegraphics[width=\linewidth]{figures/ch01/panels/fig1_2b_pc_direction.pdf}
\end{minipage}
\caption{\textbf{PCA 将最大投影方差与最小重构误差联系起来。}
同一个中心化数据集分别投影到一个非主方向和第一主成分方向上。主成分方向具有更大的投影方差，也产生更小的正交重构残差。为使图形清晰，仅对部分样本点绘出了残差线段。}
\label{fig:ch1-pca-variance-geometry}
\end{figure}

这种视角尤其重要，因为它把 PCA 重写为一对映射：
\[
x\;\longrightarrow\;z=U_k^\top x
\;\longrightarrow\;\hat{x}=U_k z.
\]
也就是说，PCA 可以被看成一个\emph{编码器}后接一个\emph{解码器}，二者都是线性的，并且都可解析求得。

\paragraph{把自编码器理解为学习得到的编码器--解码器模型。}

当 PCA 被写成重构形式后，一个自然推广就出现了。我们不再根据特征向量解析地构造编码器和解码器，而是通过最小化重构误差，直接从数据中\emph{学习}它们，由此得到自编码器的视角~\cite{baldi1989neural,hinton2006reducing}。

形式上，自编码器引入两个带参数的映射。给定输入 \(x\in\mathbb{R}^d\)：
\begin{itemize}
    \item \textbf{编码器}把输入压缩为潜在表示
    \[
    z=q_\phi(x)\in\mathbb{R}^k,\qquad k<d,
    \]
    其中 \(q_\phi\) 由参数 \(\phi\) 参数化；

    \item \textbf{解码器}根据这一潜在编码重构输入：
    \[
    \hat{x}=p_\theta(z),
    \]
    其中 \(p_\theta\) 由参数 \(\theta\) 参数化。
\end{itemize}

二者共同定义了熟悉的编码器--解码器流程
\[
x\;\longrightarrow\;z=q_\phi(x)
\;\longrightarrow\;\hat{x}=p_\theta(z).
\]
然后通过最小化重构误差来学习参数 \((\phi,\theta)\)，例如
\[
\min_{\phi,\theta}\;
\|p_\theta(q_\phi(x))-x\|^2,
\]
或对整个数据集写成
\[
\min_{\phi,\theta}\;
\sum_{i=1}^n \|p_\theta(q_\phi(x_i))-x_i\|^2.
\]

该目标与 PCA 的重构形式直接对应：
\[
\min_{U_k^\top U_k=I_k}
\sum_{i=1}^n \|U_kU_k^\top x_i-x_i\|^2.
\]
本质区别在于：PCA 中的编码器和解码器是与标准正交特征向量绑定的固定线性映射，而自编码器则把二者视为可学习函数。

\subsection{线性自编码器在什么条件下恢复 PCA？}

线性自编码器与 PCA 之间的联系是经典结果~\cite{baldi1989neural}。不过，线性自编码器本身并不能保证自动恢复 PCA。若想精确恢复 PCA，有两个结构要素十分重要。

\paragraph{1. 线性。}
假设编码器和解码器都是线性的：
\[
q_\phi(x)=W_qx,
\qquad
p_\theta(z)=W_pz,
\]
其中
\(W_q\in\mathbb{R}^{k\times d}\)、\(W_p\in\mathbb{R}^{d\times k}\)。
那么重构为
\[
\hat{x}=W_pW_qx.
\]
此时模型仍处在 PCA 的几何框架中：所有重构结果都位于 \(\mathbb{R}^d\) 中的某个 \(k\) 维线性子空间内。

如果引入非线性激活函数，这一限制就会消失。模型此时可以表示弯曲的非线性结构，而不再局限于单一平坦子空间。因此，如果目标是保持在经典 PCA 框架内，线性是必要条件。

\paragraph{2. 子空间与基的区别。}
PCA 不仅识别一个 \(k\) 维子空间，而且使用由主特征向量构成的标准正交基来描述该子空间。相比之下，线性自编码器并不会被强制学习一个标准正交基。仅仅最小化重构误差，并不会自动施加正交性约束。

这一差异很重要，但也较为细微。对于降维来说，真正关键的几何对象不是某一组特定的基，而是这些基向量所张成的\emph{子空间}。

\paragraph{子空间与列空间。}
给定向量集合 \(\{v_1,\ldots,v_k\}\subset\mathbb{R}^d\)，它们的张成空间定义为
\[
\mathrm{Span}(v_1,\ldots,v_k)
=
\left\{
\sum_{i=1}^k \alpha_i v_i:\alpha_i\in\mathbb{R}
\right\}.
\]
不同向量集合完全可能张成同一个子空间。例如
\[
\mathrm{Span}\!\left(
\begin{bmatrix}1\\0\end{bmatrix},
\begin{bmatrix}0\\1\end{bmatrix}
\right)
=
\mathrm{Span}\!\left(
\begin{bmatrix}1\\1\end{bmatrix},
\begin{bmatrix}1\\-1\end{bmatrix}
\right)
=
\mathbb{R}^2.
\]

当向量作为矩阵 \(A\in\mathbb{R}^{d\times k}\) 的列收集起来时，它们的张成空间称为 \(A\) 的\emph{列空间}（column space，也称 range）：
\[
\mathcal{R}(A)
=
\{Az:z\in\mathbb{R}^k\}.
\]

这正是降维真正关心的几何对象。PCA 使用标准正交基 \(U_k=[\,u_1,\dots,u_k\,]\) 来表示最优的 \(k\) 维子空间；线性自编码器则可能使用另一组通常并不正交的基来表示同一个子空间。

\paragraph{重构损失下的子空间恢复。}
对于线性自编码器，
\[
\hat{x}=W_pW_qx,
\]
任意重构结果都位于 \(W_p\) 的列空间中：
\[
\hat{x} \in \mathcal{R}(W_p).
\]
如果模型在平方重构误差下得到最优训练，那么这个列空间会恢复与 PCA 相同的最优 \(k\) 维子空间，即使具体的基向量并不相同。

这一点非常关键：PCA 为最优子空间选择了一组特定的标准正交基，而线性自编码器只需要恢复子空间本身即可。从重构角度看，这已经足够。只要两组基张成同一组方向，那么无论这些基是否正交，它们表示的都是同一个几何对象。

如果之后希望得到一组标准正交基，可以对学习得到的解码器矩阵进行正交化。更一般地，任何张成某个子空间的满秩向量集合，都可以使用 Gram--Schmidt 正交化、QR 分解或相关矩阵分解方法转换为同一子空间的一组标准正交基。

例如，进行 QR 分解
\[
W_p=QR,
\]
且 \(Q^\top Q=I_k\)，即可得到矩阵 \(Q\)，其各列构成与 \(W_p\) 完全相同子空间的一组标准正交基。关键在于：正交化改变的是用于表示子空间的\emph{基}，而不是子空间本身。

\subsection{从线性子空间到非线性流形}

当编码器和解码器被允许具有非线性时，自编码器真正的表达能力才显现出来。在线性情形下，如果解码器为
\[
p_\theta(z)=W_pz,
\]
那么所有重构结果都位于同一个线性子空间
\[
\{\hat{x}=W_pz:z\in\mathbb{R}^k\}
\]
中。无论潜在编码 \(z\) 如何变化，输出始终被限制在 \(\mathbb{R}^d\) 中一张平坦的 \(k\) 维“平面”上。

如果解码器变为非线性，这一几何结构将发生根本变化。所有可能重构组成的集合变为
\[
\mathcal{M}_\theta
=
\{p_\theta(z):z\in\mathbb{R}^k\},
\]
它不再被限制为一个线性子空间，而可能形成嵌入在高维环境空间中的弯曲低维曲面。这样的曲面称为\emph{流形}（manifold）。

因此，非线性自编码器以一种深刻的方式推广了 PCA。它不再只学习单一的平坦子空间，而可以学习一个弯曲流形，从而更灵活地适应数据的内在结构。

图~\ref{fig:ch1-pca-ae-manifold} 区分了两个容易混淆的概念：线性自编码器可能使用不同的基，但仍张成与 PCA 相同的子空间；而非线性解码器则真正改变了几何结构，使重构结果可以位于弯曲流形上。

\begin{figure}[H]
\centering
\begin{minipage}[t]{0.42\textwidth}
  \vspace{0pt}\centering
  \includegraphics[width=\linewidth]{figures/ch01/panels/fig1_4a_pca_basis.pdf}
\end{minipage}\hfill
\begin{minipage}[t]{0.42\textwidth}
  \vspace{0pt}\centering
  \includegraphics[width=\linewidth]{figures/ch01/panels/fig1_4b_linear_ae_basis.pdf}
\end{minipage}

\vspace{0.8em}

\begin{minipage}[t]{0.42\textwidth}
  \vspace{0pt}\centering
  \includegraphics[width=\linewidth]{figures/ch01/panels/fig1_4c_nonlinear_manifold.pdf}
\end{minipage}\hfill
\begin{minipage}[t]{0.42\textwidth}
  \vspace{0pt}
  \caption{\textbf{从 PCA 子空间到学习得到的流形。}
  PCA 使用一组特定的标准正交基表示线性子空间。线性自编码器可以用另一组通常并不正交的解码器基恢复同一个最优子空间；即使所表示的平面一致，其基向量也不必与 PCA 相同。当解码器变为非线性后，重构集合不再局限于单一平坦子空间，而可以形成弯曲的低维流形。}
  \label{fig:ch1-pca-ae-manifold}
\end{minipage}
\end{figure}

\paragraph{非线性编码器与解码器。}
在实际应用中，非线性自编码器通常由神经网络实现。编码器可以写成
\[
z=q_\phi(x),
\]
其中 \(q_\phi\) 由仿射变换与非线性激活的复合构成；解码器类似地写为
\[
\hat{x}=p_\theta(z).
\]
整体结构依然是
\[
x\;\longrightarrow\;z\;\longrightarrow\;\hat{x},
\]
但允许的变换类别已经远比 PCA 丰富。

这种从线性子空间到非线性流形的转变，是自编码器成为现代表示学习重要工具的关键原因之一。此时，PCA 可以看作一个更大范围的“基于重构的模型家族”中最简单的成员：线性、正交，并且具有解析解。

\paragraph{为什么确定性自编码器还不是生成模型。}
尽管确定性自编码器具有更强的表达能力，但它们仍缺少生成建模的一个关键要素：
\begin{quote}
\emph{对潜变量如何生成数据的概率描述。}
\end{quote}

确定性自编码器把每个输入 \(x\) 映射到一个唯一的潜在编码 \(z\)，再映射到一个唯一的重构 \(\hat{x}\)。这并没有提供有原则的方法来
\begin{itemize}
    \item 为潜在表示分配概率；
    \item 从指定分布中抽取潜变量以生成新样本；
    \item 或者依据基于似然的原则训练模型。
\end{itemize}

为了获得这些能力，我们必须超越确定性重构模型，引入\emph{概率潜变量模型}。这一转变会直接引向概率主成分分析（PPCA），随后再引向变分自编码器（VAE）；在 VAE 中，重构思想被保留下来，但被嵌入到了概率框架之中。

从这一视角看，后续路径已经十分清楚。PCA 通过方差最大化引入线性子空间；PCA 的重构形式自然引出自编码器；线性自编码器恢复与 PCA 相同的最优子空间，而非线性自编码器把这一思想推广到流形。但要真正进入现代生成建模，我们还必须再前进一步：为潜在表示赋予概率分布。

下一章将从概率主成分分析开始完成这一转变。

\section*{本章小结与参考文献}
\addcontentsline{toc}{section}{本章小结与参考文献}

本章介绍了 PCA 所需的线性代数思想，并以此建立从经典降维方法到自编码器的桥梁。

我们首先把矩阵理解为线性变换，并强调基变换、点积和投影是 PCA 背后的几何语言。由此自然引出特征向量与特征值，它们标识矩阵作用最为简单的一组特殊方向。

借助这些工具，我们先把 PCA 建立为方差最大化问题，随后又把它解释为一种将协方差矩阵对角化的基变换。这两个视角说明了为什么主成分就是最大的若干特征值对应的特征向量、为什么特征值能够度量所捕获的方差，以及 PCA 如何通过投影和重构实现降维。

随后，本章把关注点从方差转向重构。把 PCA 重新表述为线性编码器--解码器模型后，它便与自编码器建立了直接联系。线性自编码器能够恢复与 PCA 相同的最优子空间，而非线性自编码器则把这一思想从平坦子空间推广到弯曲流形。

\bigskip

本章内容建立在线性代数、矩阵分解和降维的标准理论之上。关于底层线性代数与矩阵视角，Strang~\cite{strang2016introduction} 是一部广泛适用的参考资料；Jolliffe 和 Cadima~\cite{jolliffe2016pca} 则给出了 PCA 及其解释的现代综述。就历史背景而言，PCA 的经典起源可以追溯到 Pearson~\cite{pearson1901} 和 Hotelling~\cite{hotelling1933}。Baldi 和 Hornik~\cite{baldi1989neural} 讨论了线性自编码器与 PCA 之间的联系，而 Hinton 和 Salakhutdinov~\cite{hinton2006reducing} 则对深度自编码器的发展产生了重要影响。
